\section{Related Work}
\label{sec:related-work}

\subsection{Active Inference}
\label{subsec:related-active-inference-framing}

Active inference, the framework for minimizing the free energy, is rooted in the predictive-mind view that perception is an inference process: an agent maintains beliefs about hidden causes of sensory input and updates those beliefs by reducing prediction errors~\cite{hohwy_predictive_2013}. Under the free-energy principle, perception, learning, and action can be formulated as optimization of variational free energy, a tractable objective related to Bayesian model evidence~\cite{friston_free_2006,friston_free-energy_2009,friston_action_2010,friston_predictive_2009,parr_active_2022}. 
Within active inference, a common distinction is made between variational free energy (VFE), which is used for continuous-domain perception, state estimation, and reflex-like action, and expected free energy (EFE), which is used for discrete policy selection and planning~\cite{kaplan_planning_2018}. Our contribution builds on the VFE side of active inference, more specifically on perception and self-localization, not action.

In this paper, the hidden cause of interest is the robot's field pose, and the sensory evidence consists of motion and landmark observations. The VFE formulation represents this belief in generalized coordinates, including position, velocity, and acceleration. This makes it well suited to fast-paced movement estimation, because localization is treated as a dynamically evolving belief whose higher-order motion terms directly shape prediction errors and belief updates.

\subsection{Active-Inference Robotic Navigation}
\label{subsec:related-active-inference-navigation}

Active inference has already been used in robotic navigation contexts, with a focus on operating in environments for which no prior map is available. 
Catal et al.~\cite{catal_robot_2021} connect visual perception, path integration, localization, mapping, and navigation in a hierarchical active-inference model. However, they report that ``navigation is limited to the few trajectories that were tele-operated to build the map.'' Their evaluation shows topologically consistent mapping at 200 ms observation intervals, but does not report centimeter-level metric localization accuracy. 
De Tinguy et al.~\cite{de_tinguy_online_2025} take this further by learning the environment structure online while also planning through it. Their AIMAPP framework uses motion predictions and stored panoramic observations from a 360-degree sensory view, and reports drift-robust navigation with an $(x,y)$ RMSE of $1.83 \pm 0.77$ m in a controlled warehouse. This is useful for autonomous exploration, but it is not aimed at the RoboCup requirements defined in Section~\ref{sec:motivation-background}.

Free-energy-based state estimation has also been evaluated by Bos et al.~\cite{bos_free_2021} for a quadcopter under wind-induced colored noise, showing that generalized coordinates can be used effectively in robotic state estimation. 

These active-inference papers do not address the specific challenge identified in the previous section. Our contribution is therefore not a
general navigation system, but a fast and robust localization on a known soccer field.


\subsection{RoboCup and Humanoid Localization}
\label{subsec:related-robocup}

Classical localization methods for known robot-soccer fields commonly use Kalman-filter, particle-filter, or multi-hypothesis variants to combine motion prediction with visual landmark corrections.

Jochmann et al.~\cite{jochmann_efficient_2012} are the closest classical reference for the present setting because they address robust localization in humanoid RoboCup, where repeated field landmarks can create ambiguous observations and where relocalization is important. Their multi-hypothesis unscented Kalman-filter approach is therefore a strong existing solution for known-field robot-soccer showing the capability of localization under landmark ambiguity.

However, this existing solution does not address the specific formulation pursued in this paper. It remains a classical filter-based multi-hypothesis pose estimator, rather than a variational-free-energy formulation of self-localization. It estimates pose hypotheses without generalized-coordinate dynamics, whereas fast robot-soccer play motivates a belief state in which position, velocity, and acceleration are represented together. Its UKF-style estimator uses sigma-point propagation through the landmark model, which is more computationally intensive than an explicitly derived local geometric-Jacobian update. Its motion update is driven by motor commands or control input, whereas the formulation pursued here treats odometry and IMU measurements as motion prediction errors, grounding the update in observed rather than commanded movement. Finally, because Jochmann et al. show that multi-hypothesis localization is valuable for ambiguous landmarks, their work also motivates asking how far a compact single-belief formulation can remain stable before a multi-hypothesis extension becomes necessary.


